% ==================================================================
% id:       seq-007
% topic:    סדרות
% subtopic: חישוב גבולות בעזרת אריתמטיקה של גבולות
% solution: full
% ==================================================================
\textbf{לכל אחת מהסדרות הבאות, מצאו את הגבול
$\limn a_n$:}
\begin{enumerate}\bMs{2}
    \item $\ds{\limn n (\sqrt{3n^2+1} - \sqrt{3n^2})}$

\hint{יש כאן הפרש שורשים מהצורה $\infty-\infty$: הכפילו וחלקו בצמוד.}
\begin{solbox}
    מתקיים
    \begin{align*}
        n\left(\sqrt{3n^2+1} - \sqrt{3n^2}\right)&\overset{(*)}{=}\  n\left(\sqrt{3n^2+1} - \sqrt{3n^2}\right)\cdot{\color{red}\frac{\sqrt{3n^2+1}+\sqrt{3n^2}}{\sqrt{3n^2+1}+\sqrt{3n^2}}}\\
        &=\ \frac{n}{\sqrt{3n^2+1}+\sqrt{3n^2}}\\
        &=\ \frac{n}{n\left(\sqrt{3+\frac{1}{n^2}}+\sqrt{3}\right)}\\
        &=\ \frac{1}{\sqrt{3+\frac{1}{n^2}}+\sqrt{3}}\\
    \end{align*}
    ולכן, לפי אריתמטיקה של גבולות,
    $$.\limn n\left(\sqrt{3n^2+1} - \sqrt{3n^2}\right)=\limn \frac{1}{\sqrt{3+\frac{1}{n^2}}+\sqrt{3}}=\frac{1}{2\sqrt{3}}$$
    $(\ast)$
    טריק זה נקרא כפל בצמוד
    והוא מתבסס על הזהות
    $(a-b)(a+b)=a^2-b^2$.
\end{solbox}

    \item $\ds{\limn \dfrac{n\sin(n)}{n^2+n}}$

\hint{כתבו את הסדרה כמכפלה של $\sin n$ בסדרה אחרת. מה ידוע על $\sin n$, ולאן שואפת הסדרה האחרת?}
\begin{solbox}
    נציג
    $\dfrac{n\sin(n)}{n^2-n}=a_nb_n$,
    עבור
    $$.a_n=\sin (n)\qquad, \qquad
    b_n=\dfrac{n}{n^2-n}$$
    הסדרה
    $(a_n)_{n=1}^\infty$
    הינה חסומה, אכן
    $|a_n|=|\sin (n)|\leq 1$
    לכל
    $n$.
    ואילו הסדרה
    $(b_n)_{n=1}^\infty$
    שואפת ל-
    $0$,
    אכן
    $$.\ds{\limn \dfrac{n}{n^2-n}}=\ds{\limn \dfrac{1}{n-1}}=0$$
    לכן הסדרה הנתונה שואפת ל-
    $0$
    כמכפלה של סדרה חסומה באפסה.
\end{solbox}

    \item $\ds{\limn \dfrac{n^2 + n\sin(n) - \cos(n)}{2n^2+\sqrt{n}}}$

\begin{solbox}
נחלק ב-$n^2$ כי זו החזקה הגבוהה במכנה, ונקבל
$$.\frac{n^2 + n \sin(n) - \cos(n)}{2n^2 + \sqrt{n}}
= \frac{1 + \frac{1}{n} \sin(n) - \frac{1}{n^2} \cos(n)} {2 + \frac{1}{n^{1.5}}}$$
הסדרות
$\frac{1}{n} \sin(n)$, $\frac{1}{n^2} \cos(n)$
שואפות ל-
$0$
כמכפלה של סדרה חסומה בסדרה אפסה.
גם הסדרה
$\frac{1}{n^{1.5}}$
שואפת ל-
$0$.
לכן, מאריתמטיקה של גבולות, נקבל כי
$$.\lim_{n \to \infty} \frac{1 + \frac{1}{n} \sin(n) - \frac{1}{n^2} \cos(n)} {2 + \frac{1}{n^{1.5}}}=\frac{1+0-0}{2+0}=\frac{1}{2}$$
\end{solbox}

\item $\ds{\limn \sqrt[n]{2^n + 3^n + 4^n}}$

\hint{זהו את המחובר הדומיננטי, חסמו את הביטוי מלמעלה ומלמטה בעזרתו והשתמשו בכלל הסנדוויץ'.}
\begin{solbox}
    נשים לב כי
    $4^n$
     הוא הגורם הגדול מבין שלושת המחוברים, לכן
    $$,4^n<2^n + 3^n + 4^n< 4^n+4^n+4^n=3\cdot 4^n$$
    מכאן
    $$.4=\sqrt[n]{4^n}<\sqrt[n]{2^n + 3^n + 4^n}<
    \sqrt[n]{3\cdot 4^n}=
    \sqrt[n]{3}\cdot4$$
    מכיוון ש-
    $\displaystyle\lim_{n\to\infty}\sqrt[n]{3}\cdot4=1\cdot4$,
    נסיק, מכלל הסנדוויץ', כי הגבול המבוקש הוא
    $4$.
הערה: בחישוב הגבול השתמשנו באריתמטיקה של גבולות: הגבול של מכפלה של סדרות מתכנסות שווה למכפלת הגבולות.
\end{solbox}

    \item $\ds{\limn \dfrac{2 + \sin \left(\frac{n\pi}{2}\right)}{n^{0.9}}}$
\begin{solbox}
      לכל
      $n$
      מתקיים
    $,-1\le \sin \left(\frac{n\pi}{2}\right) \le 1$
    לכן
    $$.\dfrac{2 - 1}{n^{0.9}}\leq\dfrac{2 + \sin \left(\frac{n\pi}{2}\right)}{n^{0.9}}\leq\dfrac{2 + 1}{n^{0.9}}$$
    מכיוון ש-
    $\ds\lim_{n\to\infty}\dfrac{1}{n^{0.9}}=\ds\lim_{n\to\infty}\dfrac{3}{n^{0.9}}=0$
    נסיק מכלל הסנדוויץ' כי הגבול המבוקש הוא
    $0$
    גם כן.
הערה: בחישוב הגבול השתמשנו באריתמטיקה של גבולות: הגבול של מכפלה של סדרות מתכנסות שווה למכפלת הגבולות.
\end{solbox}
    \item $\ds{\limn \dfrac{1}{n^2 + 1} + \dfrac{1}{n^2 + 2} +\cdots + \dfrac{1}{n^2 + n}}$
\hint{חסמו את הסכום מלמעלה ומלמטה על ידי החלפת כל המחוברים במחובר הגדול ביותר ובקטן ביותר, והשתמשו בכלל הסנדוויץ'.}
\begin{solbox}
       לפנינו סדרה
       $a_n$
       בה האיבר ה-$n$-י הוא סכום של
    $n$
    מחוברים. דרך אפשרית לחסום סכום שכזה היא ע"י הגדלת כל המחוברים למחובר הגדול ביותר.
    ועל מנת להקטין אותו ניתן להקטין את כל המחוברים למחובר הקטן ביותר.
    באופן כללי יותר
    $$.n\cdot\min\{x_1,x_2,\ldots,x_n\}\leq x_1+x_2+\cdots+x_n\leq n\cdot\max\{x_1,x_2,\ldots,x_n\}$$
    במקרה שלנו,
מתקיים
    $$,a_n
    =\dfrac{1}{n^2 + 1} +\cdots + \dfrac{1}{n^2 + n}\leq\underbrace{\frac{1}{n^2+1}
    +\cdots+\frac{1}{n^2+1}}_{םימעפ  \ n}=\frac{n}{n^2+1}$$
    מצד שני
    $$.a_n=\dfrac{1}{n^2 + 1} + \cdots + \dfrac{1}{n^2 + n}\geq\underbrace{\dfrac{1}{n^2 + n} + \cdots + \dfrac{1}{n^2 + n}}_{םימעפ  \ n}=\frac{1}{n+1}$$
    בסה"כ קיבלנו כי
    $$,\frac{1}{n+1}\leq a_n\leq\frac{n}{n^2+1}$$
    מכיוון ש-
    $$,\lim_{n\to\infty}\frac{n}{n^2+1}=\lim_{n\to\infty}\frac{1}{n+1}=0$$
    נסיק, ע"י שימוש בכלל הסנדוויץ', כי
    $a_n\goes 0$.
הערה: בחישוב הגבול השתמשנו באריתמטיקה של גבולות: הגבול של סכום, של מכפלה ושל מנה של סדרות מתכנסות שווה לסכום, למכפלה ולמנה של הגבולות, כל עוד הגבול של המכנה שונה מ-$0$.
\end{solbox}
    \item $\ds{\limn\frac{n}{n^2+\sin n}}$
\hint{מה הגורם הדומיננטי במונה? ובמכנה?}
\begin{solbox}
הגורם הדומיננטי במכנה הוא $n^2$. נחלק בו את המונה ואת המכנה:
$$.\frac{n}{n^2+\sin n}=\frac{\frac1n}{1+\frac{\sin n}{n^2}}$$
הסדרה $\frac{\sin n}{n^2}$ שואפת ל-$0$ כמכפלה של סדרה חסומה ($-1\le\sin n\le1$) בסדרה ששואפת ל-$0$. לכן המונה שואף ל-$0$ והמכנה שואף ל-$1$, ולפי אריתמטיקה של גבולות הגבול המבוקש הוא
    $0$.
הערה: בחישוב הגבול השתמשנו באריתמטיקה של גבולות: הגבול של סכום, של מכפלה ושל מנה של סדרות מתכנסות שווה לסכום, למכפלה ולמנה של הגבולות, כל עוד הגבול של המכנה שונה מ-$0$.
\end{solbox}

    %\item $\ds{\limn\frac{1}{\sqrt[4]{n^4+1^3+1}}+\frac{1}{\sqrt[4]{n^4+2^3+2}}+\cdots+\frac{1}{\sqrt[4]{n^4+n^3+n}}}$
\eMs
\end{enumerate}
